\documentclass[11pt]{article}

\usepackage[a4paper,margin=1in]{geometry}
\usepackage[T1]{fontenc}
\usepackage[utf8]{inputenc}
\usepackage{lmodern}
\usepackage{microtype}
\usepackage{amsmath}
\usepackage{booktabs}
\usepackage{enumitem}
\usepackage{hyperref}
\usepackage{xcolor}

\hypersetup{
  colorlinks=true,
  linkcolor=blue!50!black,
  urlcolor=blue!50!black,
  citecolor=blue!50!black
}

\title{A Fast Monte Carlo for Large Xe TPC Simulations}
\author{LXeMC Reference Manual}
\date{\today}

\begin{document}
\maketitle

\begin{abstract}
This note documents the current design of the LXeMC detector and transport
model, using the LZ detector as a concrete example of a large liquid-xenon
time projection chamber (LXe TPC). The design goal is not maximal generality.
The goal is to simulate very large event samples quickly while retaining the
physics that matters for rejection and classification. The resulting approach is
a staged simulation: a specialized fast kernel performs most of the transport
outside the fiducial volume (FV), where most events are rejected quickly, while
a full photon/electron stack is retained only inside the FV, where detailed
single-site versus multi-site behavior matters.
\end{abstract}

\section{Introduction}

Large LXe TPCs are designed to reject backgrounds aggressively. In a detector
such as LZ, most external gammas are rejected by active veto regions before they
reach the fiducial volume. This observation drives the central design choice of
the present simulation: the transport outside the FV should be optimized for
speed and early rejection, while the transport inside the FV should preserve the
detailed cascade structure needed for event topology studies.

The code therefore does not attempt to use one uniform transport model
everywhere. Instead, it exploits the detector semantics directly. Outside the
FV, the simulation only needs enough geometry and physics to decide whether a
gamma is vetoed, absorbed in passive material, or survives into the FV. Inside
the FV, the simulation keeps the full stack of photons and charged particles
because the event classification depends on the spatial pattern of energy
depositions.

\section{Problem Statement}

The practical task is to simulate a very large number of gamma events in a large
LXe detector. In the intended use case, the number of events is so large that a
uniform full simulation everywhere in the detector is too expensive. This is not
primarily a software-engineering inconvenience; it is a physics-rate problem.
The simulation must be fast enough to process event counts large enough to
support background studies at the level of interest.

The immediate consequence is that the dominant cost must be spent only on the
small subset of trajectories that can still affect the final physics observable.
Any cost spent on fully tracking obviously rejected events is wasted.

\section{Constraints and Boundary Conditions}

The design is shaped by the following constraints.

\begin{enumerate}[label=\arabic*.]
\item \textbf{Speed is dominant.} The simulation must run at the level of
millions of events per second on representative benchmark configurations.

\item \textbf{Most events are rejected early.} For many source configurations,
gammas are vetoed in active xenon outside the FV, absorbed in passive material,
or degraded to energies irrelevant for the region of interest.

\item \textbf{Detailed transport is only needed locally.} The expensive part of
the physics---full photon/electron cascades and cluster structure---matters
primarily inside the FV.

\item \textbf{The detector partition is physically meaningful.} The geometry is
not arbitrary. It should reflect the detector regions that carry distinct
transport or classification meaning.
\end{enumerate}

\section{Design Philosophy}

The design is based on a simple principle: use detector semantics to reduce the
cost of transport.

The detector is partitioned into regions that matter for either material
transport or event classification. The transport kernel is then compiled
\emph{ad hoc} from that partition. This is intentionally specialized. The
detector geometry is not interpreted as a fully generic tree at runtime in the
hot path. Instead, the canonical geometry is loaded once and transformed into a
compact fast-kernel representation with explicit scalar region parameters.

This leads to three consequences.

\begin{enumerate}[label=\arabic*.]
\item The geometry used by the fast kernel is chosen for speed rather than for
maximal abstraction.

\item The detector partition is explicit and canonical. It is not inferred
dynamically from overlapping semantic flags at runtime.

\item The full transport stack is retained only where the physics demands it,
namely inside the FV.
\end{enumerate}

\section{Simulation Strategy}

The current simulation strategy has three stages.

\subsection{Fast rejection}

The first goal is to reject events as soon as they become irrelevant. If a gamma
deposits visible energy above the veto threshold in the active xenon outside the
FV, the event is rejected immediately. If a gamma degrades below the configured
ROI floor in passive material, it is terminated locally because it can no longer
produce an FV event in the energy region of interest.

\subsection{Fast transport outside the FV}

Outside the FV, a specialized fast kernel handles gamma transport using a small
set of canonical detector regions. This kernel performs:

\begin{itemize}
\item point classification into canonical detector regions,
\item analytic distance-to-boundary calculations,
\item material-aware interaction sampling,
\item immediate veto and ROI-floor decisions.
\end{itemize}

This stage is optimized for the common case in which the gamma is rejected or
absorbed before it can matter to the final physics analysis.

\subsection{Full transport inside the FV}

If a gamma survives into the FV, the fast kernel hands off its state to a
full photon/electron transport stack restricted to the FV geometry. This stack
retains:

\begin{itemize}
\item photon propagation,
\item lepton propagation,
\item secondaries through the particle stack,
\item deposit clustering relevant to event topology.
\end{itemize}

Inside the FV, the geometry is simple: a homogeneous LXe cylinder bounded by the
fiducial radial and longitudinal cuts. The expensive part is therefore not the
geometric navigation but the cascade physics itself.

\section{Current Status and Benchmark}

The current code implements the staged architecture described above.

\begin{itemize}
\item The detector geometry is defined in two JSON files:
  \texttt{detector\_lz\_v3.json} (tracking) and
  \texttt{source\_geometry\_lz\_v1.json} (source).
\item The fast kernel is \emph{data-driven}: it reads its region list from
  JSON and identifies the tracking envelope and passive-LXe fallback by
  \texttt{role} and \texttt{tag} markers, not by hard-coded volume names.
  Adding or removing a tracked region is a JSON-only change.
\item Region capabilities (sensitive, isXe, ecut, etc.) are derived from the
  \texttt{RegionTag} enum via \texttt{\_capabilities\_for\_tag}, with optional
  per-volume JSON overrides.
\item The FV transport uses a dedicated \texttt{FVGeometry} plus a full
particle stack.
\item The legacy flat-detector workflow has been removed from the canonical
code path.
\end{itemize}

Representative benchmarks place the simulation in the target regime of roughly
$10^6$ events per second for the fast calibration workflow, with exact values
depending on source position and direction. Configurations aimed directly toward
the FV are naturally slower than configurations rejected promptly in outer veto
regions, but the code remains in the intended high-throughput regime.

\section{Geometry Model}

\subsection{Shared geometric primitives}

The file \texttt{geometry\_core.jl} contains the shared shape and wrapper layer.
It provides:

\begin{itemize}
\item primitive solids:
  \texttt{Cyl}, \texttt{CylShell}, \texttt{Box}, \texttt{Disk},
  \texttt{Cap}, \texttt{DomedContainer}, \texttt{CappedCylinder};
\item logical wrappers:
  \texttt{LCyl}, \texttt{LCylShell}, \texttt{LBox}, \texttt{LDisk};
\item materialized wrappers:
  \texttt{PhysicalVolume}, \texttt{PCyl}, \texttt{PCylShell},
  \texttt{PBox}, \texttt{PDisk};
\item common geometric operations such as containment tests, masses, and
ray entry/exit distances.
\end{itemize}

These are general geometry tools. They are not tied to the fast kernel or to the
FV stack specifically.

\subsection{Canonical detector geometry}

The file \texttt{geometry.jl} contains the canonical detector geometry
representation. It loads the tracking detector from JSON into a tree of
\texttt{TrackingDetector} nodes, with explicit region semantics and geometry.
Every volume in the JSON is required to declare a \texttt{tag} (simulation
policy) and a \texttt{role} (descriptive metadata); the code derives
per-region capabilities from the tag via \texttt{\_capabilities\_for\_tag},
with JSON fields available as overrides.

The current LZ tracking geometry is partitioned into the following regions
(defined entirely in \texttt{detector\_lz\_v3.json}, not in code):

\begin{itemize}
\item \texttt{LZ\_detector} — tracking envelope (\texttt{role: tracking\_envelope})
\item \texttt{AirDome} — gas dome above the gate
\item \texttt{LXe\_passive} — passive LXe fallback (\texttt{tag: passive\_lxe})
\item \texttt{TopActive}, \texttt{BarrelActive}, \texttt{BottomActive} — active TPC regions (\texttt{tag: tpc\_active})
\item \texttt{FV} — fiducial volume (\texttt{tag: fv})
\item \texttt{Skin} — veto layer (\texttt{tag: skin})
\item \texttt{FC\_PTFE}, \texttt{FC\_rings} — structural field-cage elements (\texttt{tag: structural})
\end{itemize}

This partition is exhaustive for the current detector model. Adding or removing
a region requires editing only the JSON file; no code changes are needed.

\subsection{Meaning of the canonical regions}

\begin{description}[leftmargin=2.0cm,style=nextline]
\item[\texttt{AirDome}] Gas region above the liquid xenon.
\item[\texttt{LXe\_passive}] Passive liquid xenon outside the active veto
regions and below the cathode.
\item[\texttt{TopActive}] Active xenon cylinder between the gate and the top of
the FV.
\item[\texttt{BarrelActive}] Active xenon shell around the FV over the FV
longitudinal span.
\item[\texttt{BottomActive}] Active xenon cylinder between the bottom of the FV
and the cathode.
\item[\texttt{FV}] Fiducial xenon volume used for event selection and topology.
\item[\texttt{Skin}] Active xenon skin outside the field-cage region.
\item[\texttt{FC\_PTFE}] PTFE part of the field cage.
\item[\texttt{FC\_rings}] Titanium ring part of the field cage.
\end{description}

\subsection{Source geometry}

The source geometry is described in a separate JSON file from the detector
tracking geometry. The separation is by \emph{contents}, not by dimensions:

\begin{itemize}
\item Source geometry hosts radioactive emission locations and source-side
      transport. It includes the cryostat Ti shells, MLI insulation, and the
      PMT structures whose effect is absorbed into the source-flux step before
      tracking begins.
\item Tracking geometry contains only what affects per-step photon transport
      after the source-flux step: the detector envelope, sensitive LXe regions,
      veto regions, and the field-cage structures that obstruct photons in
      transit.
\end{itemize}

When a volume appears in both files (\texttt{FC\_PTFE}, \texttt{FC\_rings},
\texttt{Skin}), it carries identical geometric dimensions on the two sides:
the source side is anchored to the bb0nu measured mass, and the tracking side
inherits the same R, z, and wall thickness. This means the fast-kernel
point-in-volume tests use anatomically correct dimensions and the
geometric-mass invariant is validated against the same reference on both
sides. Strict source/tracking dimensional consistency is enforced by the test
suite.

The split is therefore a \emph{subset} relation: tracking geometry is a subset
of source geometry, restricted to volumes the per-step transport actually
needs to see. It is not a license to draw the same volume with different
dimensions on the two sides.

\section{Transport Model}

\subsection{Fast kernel outside the FV}

The fast kernel operates on a compiled geometry representation generated from the
canonical detector tree. It performs:

\begin{enumerate}[label=\arabic*.]
\item region classification,
\item analytic distance-to-boundary calculation,
\item material lookup,
\item interaction-distance sampling,
\item immediate region-dependent decisions.
\end{enumerate}

This kernel treats transport in the outer detector as a problem of fast
geometric and material decisions, not as a full generic cascade simulation.

\subsection{Data-driven FastKernel construction}

The \texttt{FastKernelGeometry} is compiled from the tracking detector tree at
startup. The construction is fully data-driven:

\begin{enumerate}[label=\arabic*.]
\item Every volume in the JSON declares a \texttt{tag} (one of \texttt{world},
\texttt{vacuum}, \texttt{passive\_lxe}, \texttt{tpc\_active}, \texttt{fv},
\texttt{skin}, \texttt{structural}) and a \texttt{role} (free-form descriptive
string).

\item The function \texttt{\_capabilities\_for\_tag(tag)} provides default
capabilities for each tag: \texttt{inFastKernel}, \texttt{isXe},
\texttt{sensitive}, \texttt{ecut\_keV}, \texttt{dz\_mm}, \texttt{fv\_target}.
JSON fields can override any of these defaults per volume.

\item The compiler iterates all volumes with \texttt{inFastKernel = true} and
compiles each into a \texttt{FastKernelRegion} with precomputed scalar geometry
(rmin, rmax, zmin, zmax, cap parameters).

\item Two distinguished regions are identified by markers, not by name:
\begin{itemize}
\item \textbf{Envelope}: the unique region with \texttt{role = "tracking\_envelope"}.
\item \textbf{Fallback}: the unique region with \texttt{tag = passive\_lxe}.
\end{itemize}
The compiler validates that exactly one of each exists.

\item The \texttt{classify\_fastkernel} function tests the envelope first
(early exit if outside), then loops over all other regions in JSON declaration
order, and falls back to the passive LXe if no inner region matches.

\item The \texttt{select\_interaction\_fastkernel} function classifies
interactions by comparing integer \texttt{RegionTag} values, not volume name
strings.
\end{enumerate}

Because no volume name appears in the classification or compilation logic,
adding or removing a tracked region is a JSON-only change: add the volume entry
with the appropriate \texttt{tag}, \texttt{role}, and shape parameters, and the
FastKernel picks it up automatically.

\subsection{Full stack inside the FV}

The FV transport uses a dedicated \texttt{FVGeometry} object compiled directly
from the canonical detector geometry. The FV stack keeps:

\begin{itemize}
\item photons,
\item electrons and positrons,
\item bremsstrahlung,
\item the particle stack and deposit list.
\end{itemize}

The geometry inside the FV is reduced to what is physically needed:
the LXe material and the cylindrical FV boundary.

\subsection{Handoff logic}

The fast kernel tracks a gamma until one of the following occurs:

\begin{itemize}
\item it is vetoed outside the FV,
\item it is absorbed or terminated below the ROI floor in passive material,
\item it escapes,
\item it reaches the FV and remains relevant.
\end{itemize}

Only in the last case is the gamma handed to the FV transport stack.

\subsection{ROI floor}

The simulation uses an ROI floor in gamma energy. Once a gamma energy drops
below this threshold, the code assumes that it can no longer generate a relevant
contribution to the final ROI and terminates it locally. This is a deliberate
analysis-motivated approximation used to gain speed.

\subsection{Veto logic}

The fast kernel uses region-dependent visible-energy thresholds:

\begin{itemize}
\item active xenon outside the FV:
  veto at the configured TPC threshold;
\item skin:
  veto at the configured skin threshold;
\item passive regions:
  no veto, only transport and possible local termination.
\end{itemize}

\section{Event Logic}

\subsection{Calibration sampling}

The calibration path samples a small multiplicity of gammas and injects them
with controlled energy, position, and direction. This path is used for
benchmarks, controlled studies, and debugging.

\subsection{Veto conditions}

Deposits in active xenon outside the FV veto the event when they exceed the
configured threshold. This allows the event to be rejected immediately without
continuing detailed transport.

\subsection{FV acceptance and multi-site rejection}

Once in the FV, the event is classified using the deposit pattern returned by
the full stack. The key observable is whether the deposit pattern is consistent
with a single-site or multi-site event under the configured clustering model.

\subsection{Role of \texttt{dz\_resolution}}

The FV clustering logic uses a longitudinal clustering scale
\texttt{dz\_resolution}. Deposits separated beyond this scale contribute to
multi-site structure, while deposits consistent within this scale may be merged
into a single-site candidate.

\section{Source Flux Pipeline}

The source flux pipeline computes the rate and spectral distribution of
background gammas arriving at the detector from radioactive source volumes. It
is implemented as a three-stage process that decouples the expensive Monte Carlo
generation from the detector-side analysis.

\subsection{Factored flux tables}

The central data product of the source pipeline is a set of two-dimensional
probability tables in $(E, u)$, where $E$ is the gamma energy at the exit of the
source material and $u = \cos\theta$ is the angle of the gamma direction to the
inward surface normal (toward the detector).

For Bi-214 (single 2.448\,MeV gamma per decay), one table suffices. For Tl-208
(2.615\,MeV main gamma plus up to three companion lines at 583, 511, and
861\,keV), the main gamma and each companion line have independent tables. The
companions fire as independent Bernoulli trials with their respective branching
ratios (85\%, 23\%, 12\%), so a single decay can produce 1--4 gammas. Each
companion table also carries a survival fraction through the source material.

The tables are normalized per generated decay: $\texttt{pdf}[i,j] =
N_{\text{surviving in bin}}/N_{\text{generated}}$. Absolute rates are obtained
later by multiplying by activity $\times$ mass $\times$ BR from the decay chain.

\subsection{Compound propagation}

Gammas from the outer cryostat (OCV) must traverse the vacuum gap and then the
inner cryostat (ICV) before entering the detector. This is handled by
\texttt{propagate\_through\_layers}, which takes a sequence of
\texttt{PhysicalVolume} objects and propagates the gamma through each in order.
Between layers the gamma travels in a straight line; within each layer it
undergoes Klein-Nishina photon transport. Ray-volume entry distances use
analytic solutions for cylinders, cylindrical shells, and ellipsoidal disks.

\subsection{Cryostat source model}

For the LZ cryostat, three independent surface fluxes are computed:

\begin{description}[leftmargin=2.0cm,style=nextline]
\item[Barrel] Three contributing sources: OCV barrel $\to$ ICV
barrel (compound), MLI insulation $\to$ MLI (vacuum, straight-line) $\to$ ICV
barrel (compound), and ICV barrel self-emission. Each produces its own flux
table; an activity-weighted sum gives the combined rate.
\item[Top head] Two sources: OCV top $\to$ vacuum $\to$ ICV top (compound) and
ICV top self-emission.
\item[Bottom head] Two sources: OCV bottom $\to$ vacuum $\to$ ICV bottom
(compound) and ICV bottom self-emission.
\end{description}

Cross-contributions (e.g., OCV barrel $\to$ ICV top) are neglected.

\subsection{Virtual envelope}

Gammas sampled from the flux tables must be placed on a physical surface with a
concrete 3D position and direction before being injected into the fast kernel.
This surface is called the \emph{virtual envelope} (VE).

The VE is built from the \emph{source geometry} (\texttt{sg}), not from the
tracking detector. A geometry fix ensures that the ICV inner surface matches
the tracking detector boundary exactly: all three ICV surfaces (barrel, top
head, bottom head) share $R_{\text{inner}} = 82.1$\,cm, with head wall
thicknesses adjusted to preserve the original dome masses. This eliminates the
geometric mismatch that previously required extracting VE parameters from
\texttt{FastKernelGeometry}.

For cryostat sources the VE sits on the ICV inner surface. For PMT sources
(already inside the detector) the VE is the source volume itself. A small
inset (0.01\,cm) is applied to all VEs to avoid floating-point boundary issues
with the fast-kernel region classifier.

The \texttt{VirtualEnvelope} struct supports four surface kinds:

\begin{center}
\begin{tabular}{llrl}
\toprule
Surface & Kind & $R$ [cm] & Source \\
\midrule
Cryostat barrel & \texttt{:barrel} & 82.09 & ICV\_barrel inner \\
Cryostat top & \texttt{:cap\_up} & 82.09 & ICV\_top inner dome \\
Cryostat bottom & \texttt{:cap\_down} & 82.09 & ICV\_bottom inner dome \\
PMT top & \texttt{:disk\_flat} & 72.8 & merged PMT disk \\
PMT bottom & \texttt{:disk\_flat} & 72.8 & merged PMT disk \\
\bottomrule
\end{tabular}
\end{center}

The tracking detector geometry was also updated to match: the AirDome equator
was moved from $z = 145.6$\,cm (gate plane) to $z = 148.5$\,cm (ICV top
equator), and a new AirCyl volume fills the gas gap between the gate and the
dome.

\subsection{Three-stage background computation}

The background computation is split into three stages, each handled by a
separate script:

\begin{enumerate}[label=\arabic*.]
\item \textbf{Flux generation} (\texttt{generate\_flux\_tables.jl}). For one
(source, isotope) pair, generates $N$ decays, propagates through source
material, and writes $(E, u)$ flux tables as CSV files plus a
\texttt{metadata.json}. This is the expensive MC step; parallelized via
\texttt{julia -t N}. Run once per configuration.

\item \textbf{Background preselection} (\texttt{run\_source\_backgrounds.jl}).
Loads flux tables, samples gammas onto the virtual envelope, and processes each
event through the fast kernel and FV stack. Single-site candidates are written
to \texttt{candidates.csv} with columns $(x, y, z, E_{\text{deposited}})$.
Rejection statistics are written to \texttt{statistics.csv}. Parallelized.

\item \textbf{Final analysis} (\texttt{compute\_backgrounds.py}). Applies
Gaussian energy resolution smearing and an ROI energy window cut to the
candidates. Computes absolute background rates in counts/year. Produces summary
plots and a background table.
\end{enumerate}

This decoupling allows the flux tables to be generated once and reused with
different analysis cuts, resolution assumptions, or ROI windows without
regenerating the Monte Carlo.

\subsection{Internal sources}

Internal sources sit inside the ICV and emit directly into the LXe or gas
dome. They do not require compound propagation through cryostat layers.

\paragraph{Implemented: PMT top and bottom.}
The top PMT array (3 sub-components: PMTs, bases, structure) and bottom PMT
array (4 sub-components: PMTs, bases, structure, R8778 dome) are implemented
as merged transparent volumes. Each group is lumped into a single vacuum PDisk
spanning the full $z$ extent of its sub-components; the orientation (\texttt{:up}
for top, \texttt{:down} for bottom) controls the hemisphere filter. Since all
components share the same trivial flux table ($\sim$50\% survival from the
hemisphere cut, no material attenuation), per-component background rates are
recovered by bookkeeping: mass $\times$ specific activity.

\paragraph{Remaining: PMT barrel.}
Three transparent cylindrical-shell sources (cables, R8520 skin PMTs,
R8778 lower-ring PMTs) at $R \approx 81.2$--$81.4$\,cm along the FC outer
wall. Same pattern as PMT top/bottom but with a barrel VE.

\paragraph{Remaining: field cage sources.}
Six sources inside or on the field cage: FC\_PTFE and FC\_rings (dense, KN
self-shielding), FC\_topgrid and FC\_botgrid (stainless steel, KN),
FC\_resistors and FC\_sensors (transparent). These emit directly into the LXe
with single-source flux generation (no compound propagation). Each source has
its own dispatch entry and barrel-type VE.

\section{Code Map}

The current code ownership is as follows.

\begin{description}[leftmargin=2.5cm,style=nextline]
\item[\texttt{geometry\_core.jl}] shared geometric primitives, wrappers,
ray entry/exit (analytic for Cyl, CylShell, Disk), and geometric helpers.
\item[\texttt{geometry.jl}] canonical detector geometry, tracking tree, region
semantics, fast-kernel compilation, FV geometry compilation.
\item[\texttt{tracking.jl}] fast transport outside the FV, FV-only full stack,
event-level logic. \texttt{process\_event} returns FV deposits for accepted
events.
\item[\texttt{source\_propagation.jl}] photon-only KN transport inside source
volumes (with vacuum shortcut for transparent sources), uniform position
sampling, direction classification.
\item[\texttt{source\_flux.jl}] flux table structs (\texttt{SourceFluxBi214},
\texttt{SourceFluxTl208}, \texttt{SourceRateTable}), source geometry loader,
single-source and compound flux generators.
\item[\texttt{source\_sampling.jl}] 2D PDF sampling, surface point generation
(barrel, ellipsoidal caps, flat disks), direction reconstruction,
gamma-from-flux sampling.
\item[\texttt{source\_dispatch.jl}] \texttt{VirtualEnvelope} (built from
source geometry for all sources), \texttt{make\_surface\_sampler}, source
dispatch (\texttt{dispatch\_source\_flux}) and thread-result merging.
\item[\texttt{cryostat\_sources.jl}] cryostat-specific flux functions (barrel,
top, bottom) with activity-weighted rate table construction.
\item[\texttt{pmt\_sources.jl}] PMT flux functions (\texttt{pmt\_top\_flux},
\texttt{pmt\_bottom\_flux}), merged transparent volume construction with
per-component rate bookkeeping.
\item[\texttt{flux\_utils.jl}] flux table merging, CSV I/O, JSON metadata
helpers.
\item[\texttt{decays.jl}] decay and calibration sampling.
\end{description}

The public entry points of primary interest are:

\begin{itemize}
\item \texttt{load\_tracking\_detector},
      \texttt{compile\_fastkernel\_geometry},
      \texttt{compile\_fv\_geometry}
\item \texttt{process\_event},
      \texttt{propagate\_gamma\_in\_fv}
\item \texttt{load\_source\_geometry},
      \texttt{make\_virtual\_envelope},
      \texttt{make\_surface\_sampler}
\item \texttt{dispatch\_source\_flux},
      \texttt{supported\_sources},
      \texttt{supported\_isotopes}
\item \texttt{generate\_flux\_bi214},
      \texttt{generate\_flux\_tl208},
      \texttt{generate\_flux\_compound\_bi214},
      \texttt{generate\_flux\_compound\_tl208}
\item \texttt{cryostat\_barrel\_flux},
      \texttt{cryostat\_top\_flux},
      \texttt{cryostat\_bottom\_flux}
\item \texttt{pmt\_top\_flux},
      \texttt{pmt\_bottom\_flux}
\item \texttt{sample\_from\_flux},
      \texttt{sample\_gamma\_from\_flux}
\end{itemize}

\section{Current Limitations}

The current simulation is fast and purposeful, but it is not intended to be an
arbitrary general detector framework. The main limitations are:

\begin{enumerate}[label=\arabic*.]
\item The fast kernel reads its region list from JSON and is data-driven, but
it assumes coaxial cylindrical geometry (all volumes centered on the detector
axis). Non-coaxial or non-cylindrical detector topologies would require
extending the kernel's geometric predicates.
\item The ROI-floor termination is an intentional approximation.
\item The outside-FV transport is optimized for rapid rejection rather than for
full generic cascade tracking in every material.
\item The \texttt{RegionTag} enum defines a fixed set of simulation policies
(\texttt{TAG\_FV}, \texttt{TAG\_TPC\_ACTIVE}, \texttt{TAG\_SKIN}, etc.). A
detector with a genuinely new kind of sensitive region would require extending
the enum and \texttt{\_capabilities\_for\_tag}.
\end{enumerate}

These limitations are acceptable because they follow directly from the project
goal: simulate extremely large event samples as fast as possible while retaining
the detailed physics only where it matters.

\end{document}
